% =====================================================================
%  FICHE 8 — THÈME 4 — NIVEAU MOYEN — CORRIGÉ
% =====================================================================
\documentclass[11pt,a4paper]{article}
\usepackage[utf8]{inputenc}
\usepackage[T1]{fontenc}
\usepackage{lmodern}
\usepackage[french]{babel}
\usepackage{amsmath,amssymb,mathtools}
\usepackage{icomma}
\usepackage{siunitx}
\sisetup{output-decimal-marker={,},per-mode=power,group-separator={\,},group-minimum-digits=5}
\usepackage[version=4]{mhchem}
\usepackage{xcolor}
\usepackage{tikz}
\usepackage[most]{tcolorbox}
\usepackage{enumitem}
\usepackage{array}
\usepackage{fancyhdr}
\usepackage[hidelinks]{hyperref}
\usetikzlibrary{arrows.meta,positioning,calc}
\usepackage[a4paper,margin=2.1cm,headheight=26pt,headsep=10pt]{geometry}
\usepackage{titlesec}

\definecolor{primary}{RGB}{150,98,162}
\definecolor{primarydark}{RGB}{92,52,110}
\definecolor{excol}{RGB}{0,135,90}

\tcbset{boxbase/.style={enhanced,breakable,boxrule=0.9pt,arc=2.5pt,left=3mm,right=3mm,top=2.3mm,bottom=2.3mm,fonttitle=\bfseries,coltitle=white,toptitle=0.6mm,bottomtitle=0.6mm}}
\newtcolorbox[auto counter]{correction}[1][]{boxbase,colback=excol!4,colframe=excol,title={\textsc{Corrigé}~\thetcbcounter\ \ifx\relax#1\relax\relax\else\textmd{--- \itshape #1}\fi}}

\newcommand{\dS}{\Delta S}
\newcommand{\rep}[1]{\textcolor{primarydark}{\textbf{#1}}}

\pagestyle{fancy}
\fancyhf{}
\fancyhead[L]{\small\textcolor{primarydark}{\textbf{Fiche 8} — Thème 4 : Entropie et deuxième principe}}
\fancyhead[R]{\small\textcolor{primarydark}{\textsc{Moyen — corrigé}}}
\fancyfoot[C]{\small\thepage}
\renewcommand{\headrulewidth}{0.8pt}
\renewcommand{\headrule}{\hbox to\headwidth{\color{primary}\leaders\hrule height \headrulewidth\hfill}}
\setlength{\parindent}{0pt}\setlength{\parskip}{3pt}

\begin{document}

\begin{center}
\begin{tikzpicture}
\node[rectangle,rounded corners=7pt,fill=excol,text=white,minimum width=16cm,
      minimum height=1.1cm,align=center,inner sep=6pt]
  {{\large\bfseries Thème 4 — Entropie et deuxième principe}\\[1pt]{\small Niveau \textsc{Moyen} — corrigé détaillé}};
\end{tikzpicture}
\end{center}

\begin{correction}[compter les moles de gaz]
\begin{enumerate}[label=\alph*),leftmargin=2em,topsep=2pt,itemsep=2pt]
  \item $3$ mol de gaz $\to 2$ mol : moins de gaz, \rep{$\dS<0$}.
  \item $0 \to 1$ mol de gaz : apparition de gaz, \rep{$\dS>0$}.
  \item $2$ mol de gaz $\to 0$ (solide) : \rep{$\dS<0$}.
  \item $0 \to 1$ mol de gaz ($\ce{O2}$), les liquides comptent peu : \rep{$\dS>0$}.
  \item $2$ mol de gaz $\to 2$ mol : nombre de moles de gaz inchangé, \rep{$\dS\approx0$}.
\end{enumerate}
\end{correction}

\begin{correction}[le gaz est le facteur décisif]
On ne compte que les \emph{moles de gaz}.
\begin{enumerate}[label=\alph*),leftmargin=2em,topsep=2pt,itemsep=2pt]
  \item $1$ mol de gaz $\to 1$ mol : \rep{$\dS\approx0$} (le solide qui disparaît ne change pas le compte de gaz).
  \item $1$ mol de gaz $\to 2$ mol : \rep{$\dS>0$}.
  \item $3$ mol de gaz $\to 3$ mol : \rep{$\dS\approx0$}.
  \item $3$ mol de gaz $\to 0$ (tout devient solide) : \rep{$\dS<0$}.
\end{enumerate}
\end{correction}

\begin{correction}[entropie et réaction réversible]
\begin{enumerate}[label=\alph*),leftmargin=2em,topsep=2pt,itemsep=2pt]
  \item Sens direct : $1$ mol de gaz $\to 2$ mol, le désordre augmente, \rep{$\dS>0$}.
  \item Le désordre favorise le sens qui \emph{crée} le plus de gaz : le \rep{sens direct}
        ($\ce{N2O4}\to2\,\ce{NO2}$).
  \item Sens inverse : \rep{$\dS<0$} (on repasse de $2$ à $1$ mol de gaz).
\end{enumerate}
\end{correction}

\begin{correction}[désordre et spontanéité]
\begin{enumerate}[label=\alph*),leftmargin=2em,topsep=2pt,itemsep=2pt]
  \item Solide $\to$ liquide : le désordre augmente, \rep{$\dS_{\text{système}}>0$}.
  \item La fusion est spontanée à \SI{25}{\celsius}, donc d'après le deuxième principe
        \rep{$\dS_{\text{univers}}>0$}.
  \item Le glaçon fond \rep{entièrement} : le processus va jusqu'au bout (l'eau liquide, plus désordonnée,
        est l'état stable à \SI{25}{\celsius}).
\end{enumerate}
\end{correction}

\end{document}
